\chapter{Das Geschlecht}

\section{Eine Zahl, die alles entscheidet}

\begin{merksatz}
\boxtitel{Hilbert-Polynome projektiver Kurven}
Für eine projektive Kurve zählt man homogene Polynome vom Grad genau $d$ modulo dem Ideal --- die projektive Hilbert-Funktion. Gerechnet:

\begin{center}
{\small
\renewcommand{\arraystretch}{1.46}
\begin{tabular}{@{}P{3.4cm}P{2.6cm}P{1.0cm}P{1.4cm}@{}}
\toprule
\textbf{Kurve} & $H(d)$ für $d\ge d_0$ & \textbf{Grad} & $g$\\
\midrule
Gerade & $d + 1$ & $1$ & $0$\\
Kegelschnitt $x^{2} + y^{2} = z^{2}$ & $2d + 1$ & $2$ & $0$\\
verdrehte Kubik (im $\mathbb P^{3}$) & $3d + 1$ & $3$ & $0$\\
glatte Kubik $y^{2}z = x^{3} - 2xz^{2}$ & $3d$ & $3$ & $1$\\
Kubik mit Knoten $y^{2}z = x^{3} + x^{2}z$ & $3d$ & $3$ & $1$ (!)\\
glatte Quartik $x^{4} + y^{4} = z^{4}$ & $4d - 2$ & $4$ & $3$\\
\bottomrule
\end{tabular}}
\end{center}

\begin{center}
\begin{minipage}{0.9\textwidth}
\textbf{Das Hilbert-Polynom einer Kurve ist $ed + 1 - g$: Der Leitkoeffizient ist der Grad, die Konstante bestimmt eine Zahl $g$ --- das Geschlecht. Für alle Kurven, die Kapitel 11 parametrisiert hat, ist $g = 0$. Für die glatte Kubik ist $g = 1$. Die Kubik mit Knoten hat dasselbe Polynom wie die glatte --- und ist rational. Das Ausrufezeichen ist der Bruchfall dieses Kapitels.}
\end{minipage}
\end{center}

\begin{center}
\begin{minipage}{0.88\textwidth}
\kursiv{\textbf{Das Geschlecht ist die wichtigste Zahl der Kurventheorie}: Es entscheidet, ob eine Kurve rational ist ($g = 0$), ob sie ein Gruppengesetz trägt ($g = 1$), und ob sie nur endlich viele rationale Punkte hat ($g\ge2$, Faltings). Es ist topologisch (die Zahl der Henkel der komplexen Kurve als Fläche), algebraisch (die Dimension des Raums der Differentialformen), kombinatorisch (aus Grad und Singularitäten berechenbar) --- und alle Definitionen stimmen überein, mit einer Ausnahme, die man kennen muss. Dieses Kapitel gibt die berechenbare Fassung, beweist sie über das Hilbert-Polynom und zieht die Folgerungen, die der Band versprochen hat: Fermat, Mordell, Faltings.}
\end{minipage}
\end{center}
\end{merksatz}

\section{Grad, Geschlecht, Singularitäten}

\begin{definition}
\boxtitelm{Definition 12.1 --- arithmetisches Geschlecht}
Für eine projektive Kurve $C\subseteq\mathbb P^{n}$ mit Hilbert-Polynom $HP_C(d) = ed + c$ heißt $e = \deg C$ der Grad und $p_a(C) = 1 - c$ das \Begriff{arithmetische Geschlecht}. Für eine glatte Kurve heißt $g = p_a$ schlicht das \Begriff{Geschlecht}.
\end{definition}

\begin{merksatz}
\boxtitel{Satz 12.2 --- die Grad-Geschlecht-Formel}
Eine ebene projektive Kurve vom Grad $d$ hat arithmetisches Geschlecht
\[ p_a = \frac{(d - 1)(d - 2)}{2} . \]
Ist sie glatt, so ist das ihr Geschlecht: Geraden und Kegelschnitte $0$, Kubiken $1$, Quartiken $3$, Quintiken $6$. Hat sie nur gewöhnliche Doppelpunkte (Knoten) als Singularitäten, $\delta$ an der Zahl, so ist ihr \Begriff{geometrisches Geschlecht} --- das Geschlecht der Kurve nach Auflösung der Singularitäten ---
\[ g = \frac{(d - 1)(d - 2)}{2} - \delta . \]

\schritt{Beweis der ersten Formel}
Sei $C = V(F)$, $F$ homogen vom Grad $d$ in $x, y, z$. Die homogenen Polynome vom Grad $m$ bilden einen Raum der Dimension $\binom{m + 2}{2}$. Die Elemente des Ideals $\langle F\rangle$ vom Grad $m$ sind $F\cdot G$ mit $\deg G = m - d$: Dimension $\binom{m - d + 2}{2}$ für $m\ge d$. Also
\begin{align*}
H_C(m) &= \binom{m + 2}{2} - \binom{m - d + 2}{2} = \frac{(m + 2)(m + 1) - (m - d + 2)(m - d + 1)}{2}\\
&= dm + \frac{3d - d^{2}}{2},
\end{align*}
nach Ausmultiplizieren. Leitkoeffizient $d$ \checkmark; Konstante $c = \frac{3d - d^{2}}{2}$, also $p_a = 1 - c = \frac{d^{2} - 3d + 2}{2} = \frac{(d - 1)(d - 2)}{2}$. $\blacksquare$

\schritt{Probe}
$d = 3$: $H(m) = 3m + 0$ \checkmark\ (Tabelle). $d = 4$: $4m - 2$ \checkmark. $d = 2$: $2m + 1$ \checkmark.

\schritt{Zur zweiten Formel}
Jeder Knoten senkt das Geschlecht um $1$: Die Auflösung trennt die zwei Zweige, und die Kurve verliert einen \enquote{Henkel}. Der Beweis geht über Differentialformen oder über das Hilbert-Polynom der Normalisierung (Auslassung 17, \zit{7}, Kap.~8; \zit{10}, Kap.~IV.3). Für eine Spitze gilt dasselbe ($\delta = 1$); für kompliziertere Singularitäten ist $\delta$ eine lokale Invariante, die größer als $1$ sein kann.

\begin{center}
{\small
\renewcommand{\arraystretch}{1.5}
\begin{tabular}{@{}P{3.4cm}P{0.8cm}P{1.0cm}P{1.0cm}P{2.0cm}@{}}
\toprule
\textbf{Kurve} & $d$ & $\delta$ & $g$ & \textbf{Kapitel 11}\\
\midrule
Kreis & $2$ & $0$ & $0$ & rational\\
Kubik mit Knoten / Spitze & $3$ & $1$ & $0$ & rational\\
Folium & $3$ & $1$ & $0$ & rational\\
Lemniskate & $4$ & $3$ & $0$ & rational (ein reeller Knoten, zwei in den Kreispunkten)\\
glatte Kubik & $3$ & $0$ & $1$ & elliptisch\\
glatte Quartik & $4$ & $0$ & $3$ & ---\\
Leitfall-Hyperbel & $2$ & $0$ & $0$ & rational\\
\bottomrule
\end{tabular}}
\end{center}

\begin{center}
\begin{minipage}{0.88\textwidth}
\kursiv{\textbf{Eine irreduzible Kurve kann höchstens $\binom{d-1}{2}$ Singularitäten haben} (das Geschlecht ist nie negativ) --- das ist die Aussage aus Aufgabe 8.1, und bei Gleichheit ist sie rational. Die Lemniskate schöpft die Schranke aus: $\binom32 = 3$ Knoten, Geschlecht null, rational. Hätte sie nur den einen sichtbaren Knoten, wäre sie vom Geschlecht $2$ und nicht parametrisierbar --- die zwei unsichtbaren Doppelpunkte in den Kreispunkten sind entscheidend, und man sieht sie nur projektiv und komplex. Das ist der Grund, warum Kapitel 9 vor Kapitel 12 kommt.}
\end{minipage}
\end{center}
\end{merksatz}

\section{Was das Geschlecht entscheidet}

\begin{merksatz}
\boxtitel{Satz 12.3 --- die Trichotomie}
Für eine glatte projektive Kurve $C$ über $\mathbb C$:
\begin{center}
{\small
\renewcommand{\arraystretch}{1.5}
\begin{tabular}{@{}P{1.4cm}P{3.0cm}P{4.2cm}@{}}
\toprule
$g$ & \textbf{Geometrie} & \textbf{Arithmetik (über $\mathbb Q$)}\\
\midrule
$0$ & $C\cong\mathbb P^{1}$: rational & keine oder unendlich viele rationale Punkte, parametrisierbar\\
$1$ & elliptisch: ein Torus & endlich erzeugte Gruppe (Mordell), Rang unbekannt\\
$\ge2$ & \enquote{allgemeiner Typ}: Fläche mit $g$ Henkeln & endlich viele rationale Punkte (Mordell-Vermutung, Faltings 1983)\\
\bottomrule
\end{tabular}}
\end{center}

\schritt{Beweis für $g = 0$}
Eine glatte Kurve vom Geschlecht null über $\mathbb C$ ist isomorph zu $\mathbb P^{1}$: Das ist Riemann--Roch (Auslassung 18) --- mit $g = 0$ gibt es zu jedem Punkt $P$ eine rationale Funktion mit genau einem einfachen Pol in $P$, und sie ist ein Isomorphismus nach $\mathbb P^{1}$. Über $\mathbb Q$ braucht es zusätzlich einen rationalen Punkt (Kapitel 11). $\blacksquare$

\schritt{Zu $g = 1$ und $g\ge2$}
Mordell: Kapitel 10, Auslassung 15. Faltings' Satz ist einer der tiefsten der Mathematik des 20.\ Jahrhunderts, und sein Beweis liegt weit außerhalb dieses Bandes (Auslassung 19); die Aussage aber ist einfach: \textbf{Eine Kurve vom Geschlecht $\ge2$ hat nur endlich viele rationale Punkte.}

\schritt{Fermat}
$x^{n} + y^{n} = z^{n}$ ist eine glatte Kurve vom Grad $n$ (die partiellen Ableitungen $nx^{n-1}, ny^{n-1}, nz^{n-1}$ verschwinden nur im Nullpunkt), also vom Geschlecht $\frac{(n-1)(n-2)}{2}$: $0$ für $n = 2$ (die pythagoreischen Tripel, unendlich viele), $1$ für $n = 3$ (Eulers Fall, Rang $0$), $\ge3$ für $n\ge4$. Faltings' Satz gibt also für jedes $n\ge4$: nur endlich viele primitive Lösungen. Dass es gar keine sind, ist Wiles 1995 --- über die Modularität elliptischer Kurven, Kapitel 10 --- und nicht aus dem Geschlecht allein.

\begin{center}
\begin{minipage}{0.9\textwidth}
\textbf{Das Geschlecht teilt alle Kurven in drei Welten: die rationale, in der man parametrisiert; die elliptische, in der man addiert; und die allgemeine, in der die rationalen Punkte endlich sind und niemand sie systematisch finden kann.}
\end{minipage}
\end{center}

\begin{center}
\begin{minipage}{0.88\textwidth}
\kursiv{\textbf{Die Trichotomie ist das Muster, das sich durch die ganze Geometrie zieht}: Riemanns Uniformisierung ordnet den drei Geschlechtern die drei Geometrien zu --- sphärisch ($g = 0$, positive Krümmung), euklidisch ($g = 1$, flach, der Torus), hyperbolisch ($g\ge2$, negative Krümmung) ---, und Langs Vermutung überträgt sie auf höhere Dimensionen: Varietäten \enquote{vom allgemeinen Typ} sollen nur wenige rationale Punkte haben. Für den Leitfall ist alles Geschlecht null: Die Fasern linearer Kompartimentmodelle sind rationale Varietäten, weil die Messgleichungen in jedem Parameter linear oder quadratisch sind. Nichtlineare Modelle --- Michaelis--Menten-Kinetik, Hill-Funktionen --- können Fasern höheren Geschlechts haben, und dann gibt es keine Parametrisierung der nicht identifizierbaren Parameter; Kapitel 13.}
\end{minipage}
\end{center}
\end{merksatz}

\section{Der Bruchfall}

\begin{warnung}
\boxtitelr{Zwei Geschlechter}
Die Kubik mit Knoten hat das Hilbert-Polynom $3d$ --- arithmetisches Geschlecht $1$ --- und ist rational --- geometrisches Geschlecht $0$. Beide Zahlen sind richtig; sie messen Verschiedenes.

\begin{center}
{\small
\renewcommand{\arraystretch}{1.5}
\begin{tabular}{@{}P{3.0cm}P{5.8cm}@{}}
\toprule
\textbf{Arithmetisches Geschlecht $p_a$} & aus dem Hilbert-Polynom; sieht die Kurve, wie sie im $\mathbb P^{n}$ liegt, mit Singularitäten; $p_a = \binom{d-1}{2}$ für jede ebene Kurve vom Grad $d$\\
\textbf{Geometrisches Geschlecht $g$} & Geschlecht der Normalisierung (Singularitäten aufgelöst); $g = p_a - \delta$; nur $g$ entscheidet über Rationalität\\
\textbf{Gleichheit} & genau für glatte Kurven\\
\bottomrule
\end{tabular}}
\end{center}

\begin{center}
\begin{minipage}{0.9\textwidth}
\textbf{Wer das Geschlecht aus dem Hilbert-Polynom liest und nicht nach Singularitäten sucht, erklärt rationale Kurven für elliptisch. Die Rechnung ist richtig, die Zahl ist die falsche.}
\end{minipage}
\end{center}

\begin{center}
\begin{minipage}{0.88\textwidth}
\kursiv{\textbf{Auch über $\mathbb R$ versagt die Anschauung}: Die reelle Kubik $y^{2} = x^{3} - 2x$ hat zwei Zusammenhangskomponenten (Titelbild), $y^{2} = x^{3} + x$ nur eine --- beide haben Geschlecht $1$, und die komplexe Kurve ist in beiden Fällen ein Torus. Die Zahl der reellen Komponenten einer Kurve vom Geschlecht $g$ ist höchstens $g + 1$ (Harnack 1876), und welche Zahlen zwischen $0$ und $g + 1$ vorkommen, ist für Kurven vollständig bekannt, für Flächen nicht. Das Geschlecht ist eine komplexe Invariante; was man reell sieht, ist eine ihrer vielen Erscheinungsformen.}
\end{minipage}
\end{center}
\end{warnung}

\section{Übungsaufgaben}

\begin{aufgabe}
\boxtitel{Aufgabe 12.1 --- Geschlecht einer Raumkurve}
Die verdrehte Kubik hat Hilbert-Polynom $3d + 1$. Geschlecht? Und warum stimmt die Grad-Geschlecht-Formel hier nicht?

\loesung
\schritt{Schritt 1}
$3d + 1 = ed + 1 - g$ mit $e = 3$: $g = 0$. Rational --- stimmt, $(t, t^{2}, t^{3})$.

\schritt{Schritt 2 --- Die Formel}
$\binom{d-1}{2} = \binom22 = 1$ für $d = 3$ --- aber das gilt für \kursiv{ebene} Kurven. Die verdrehte Kubik liegt im $\mathbb P^{3}$, nicht in einer Ebene; projiziert man sie in eine Ebene, erhält man eine ebene Kubik mit einem Knoten (oder einer Spitze), Geschlecht $1 - 1 = 0$. \checkmark

\schritt{Schritt 3 --- Allgemein}
Für Kurven im $\mathbb P^{3}$ vom Grad $d$ ist das Geschlecht höchstens $\binom{d-1}{2}$ (ebene Kurven haben das größte), und Castelnuovo gab 1889 die scharfe Schranke für Kurven, die in keiner Ebene liegen: etwa $\frac{d^{2}}{4} - d + 1$. Für $d = 3$: $0$. \checkmark\ Raumkurven haben kleineres Geschlecht als ebene gleichen Grades.

\begin{center}
\begin{minipage}{0.88\textwidth}
\kursiv{\textbf{Die Grad-Geschlecht-Formel ist eine Aussage über die Einbettung, nicht über die Kurve}: Dieselbe abstrakte Kurve --- $\mathbb P^{1}$ --- erscheint als Gerade (Grad $1$), als Kegelschnitt (Grad $2$), als verdrehte Kubik (Grad $3$), als rationale Normkurve beliebigen Grades, und jedes Mal mit Geschlecht $0$. Der Grad ist eine Eigenschaft der Lage im Raum, das Geschlecht eine der Kurve selbst. Das Hilbert-Polynom enthält beides: den Grad vorn, das Geschlecht hinten.}
\end{minipage}
\end{center}
\end{aufgabe}

\begin{aufgabe}
\boxtitel{Aufgabe 12.2 --- Hyperelliptische Kurven}
$y^{2} = f(x)$ mit $f$ vom Grad $5$ oder $6$ ohne mehrfache Nullstellen. Geschlecht? (Vorsicht im Unendlichen.)

\loesung
\schritt{Schritt 1 --- Die naive Rechnung}
Grad $6$ in der Ebene: $\binom52 = 10$. Das kann nicht stimmen --- die Kurve ist \enquote{zweiblättrig} über der $x$-Achse und sollte einfacher sein.

\schritt{Schritt 2 --- Das Unendliche}
Homogen: $y^{2}z^{4} = F(x, z)$ mit $F$ vom Grad $6$. Bei $z = 0$: $x^{6} = 0$, ein Punkt $(0 : 1 : 0)$, und dort ist die Kurve \kursiv{hochgradig singulär} (Vielfachheit $4$: der Term $y^{2}z^{4}$ hat in den lokalen Koordinaten $x/y, z/y$ den Grad $4$). Der Beitrag $\delta$ dieser Singularität ist $\binom{4}{2}\cdot\ldots$ --- genauer: Die Auflösung zeigt $\delta = 8$ für $\deg f = 6$.

\schritt{Schritt 3 --- Das Geschlecht}
$g = 10 - 8 = 2$. Für $\deg f = 5$: ebene Quintik $\binom42 = 6$, Singularität im Unendlichen mit $\delta = 4$, $g = 2$. Allgemein: $y^{2} = f(x)$ mit $\deg f = 2g + 1$ oder $2g + 2$ hat Geschlecht $g$ --- die \Begriff{hyperelliptischen Kurven}.

\schritt{Schritt 4 --- Richtig gerechnet}
Besser als über die singuläre ebene Form: Riemann--Hurwitz. Die Projektion $(x, y)\mapsto x$ ist zweiblättrig über $\mathbb P^{1}$ mit Verzweigung in den $2g + 2$ Nullstellen von $f$ (und für ungeraden Grad zusätzlich im Unendlichen): $2g - 2 = 2\cdot(0 - 2) + (2g + 2)$, also $g = g$. \checkmark\ (Auslassung 18.)

\begin{center}
\begin{minipage}{0.88\textwidth}
\kursiv{\textbf{Für $g = 2$ hat die Kurve nur endlich viele rationale Punkte} (Faltings), und seit \Person{Chabauty} 1941 und \Person{Coleman} 1985 gibt es ein Verfahren, sie in vielen Fällen zu finden --- für $y^{2} = x^{5} - x + 1$ oder die Kurven, die bei Diophant stehen. \Person{Wetherell} löste 1997 so ein Problem aus Diophants Arithmetica VI.17, das 1700 Jahre offen war. Für den Leitfall: Eine Faser vom Geschlecht $2$ hieße, dass die nicht identifizierbaren Parameter nicht einmal durch eine rationale Funktion parametrisiert werden könnten --- ein Fall, der bei nichtlinearen Reaktionsnetzen vorkommt und dort die Modellreduktion unmöglich macht.}
\end{minipage}
\end{center}
\end{aufgabe}

\begin{aufgabe}
\boxtitel{Aufgabe 12.3 --- Das Geschlecht der Leitfall-Faser}
Für ein Modell mit zwei freien Parametern sei die Faser die Kurve $u^{2}v + uv^{2} = u + v$ in den nicht identifizierbaren Parametern $u, v$. Geschlecht, Rationalität, und was folgt für die Modellreduktion?

\loesung
\schritt{Schritt 1 --- Grad und Singularitäten}
$f = u^{2}v + uv^{2} - u - v = uv(u + v) - (u + v) = (u + v)(uv - 1)$: \textbf{reduzibel}! Zwei Komponenten: die Gerade $u + v = 0$ und die Hyperbel $uv = 1$.

\schritt{Schritt 2 --- Geschlecht je Komponente}
Beide vom Geschlecht $0$: Gerade, Kegelschnitt. Jede ist rational.

\schritt{Schritt 3 --- Modellreduktion}
Die nicht identifizierbaren Parameter liegen entweder auf $v = -u$ oder auf $v = 1/u$. Auf der ersten Komponente ist $u + v$ identifizierbar (null), auf der zweiten $uv$ (eins) --- verschiedene identifizierbare Kombinationen auf verschiedenen Komponenten. Ein Modell, dessen Faser reduzibel ist, hat keine einheitliche Reduktion: Man muss erst wissen, auf welcher Komponente man ist.

\schritt{Schritt 4 --- Hätte man das Geschlecht $\binom{2}{2} = 1$ gelesen}
$d = 3$, $p_a = 1$, und man hätte \enquote{elliptische Faser, keine Parametrisierung} geschlossen --- falsch, die Kurve ist reduzibel, und $p_a$ einer reduziblen Kurve ist nicht das Geschlecht eines Bestandteils (für zwei Komponenten vom Geschlecht $0$, die sich in $k$ Punkten schneiden, ist $p_a = k - 1$; hier $k = 2$: $p_a = 1$. \checkmark).

\begin{center}
\begin{minipage}{0.88\textwidth}
\kursiv{\textbf{Zuerst zerlegen, dann das Geschlecht} --- der Bruchfall von Kapitel 7 in neuer Gestalt. Für die Praxis der Identifizierbarkeit heißt das: Die Faser über generischen Daten wird in irreduzible Komponenten zerlegt (Kapitel 3, mit Gröbner-Basen), jede Komponente bekommt Dimension, Grad und --- wenn sie eine Kurve ist --- Geschlecht; und erst dann weiß man, ob die nicht identifizierbaren Parameter sich durch eine rationale Funktion ausdrücken lassen. Die Programme tun das noch nicht; es ist eine der offenen Aufgaben, Kapitel 13.}
\end{minipage}
\end{center}
\end{aufgabe}

\begin{geschichte}
\boxtitelg{Abel, Riemann, Clebsch}
\Person{Niels Henrik Abel} fand 1826, dass die Integrale $\int\mathrm dx/\sqrt{f(x)}$ für $\deg f = 2g + 1$ oder $2g + 2$ ein Additionstheorem mit $g$ Summanden haben --- die Zahl $g$ erschien zum ersten Mal, als Zahl der Integrale, die man braucht. \Person{Riemann} erkannte 1857 in $g$ die Zahl der Henkel der zugehörigen Fläche und nannte sie --- oder \Person{Clebsch} 1865 --- das Geschlecht.

\begin{center}
\begin{minipage}{0.9\textwidth}
\kursiv{\textbf{Die Grad-Geschlecht-Formel ist von \Person{Clebsch} 1864}, die Fassung mit Singularitäten von \Person{Max Noether} und \Person{Brill} 1874; \Person{Hilbert} 1890 lieferte das Polynom, aus dem man beides abliest, und \Person{Severi} nannte 1909 die Konstante das arithmetische Geschlecht. Dass das Geschlecht alles entscheidet, vermutete \Person{Mordell} 1922 in derselben Arbeit, in der er die endliche Erzeugung bewies: \enquote{Eine Kurve vom Geschlecht $\ge2$ hat nur endlich viele rationale Punkte} --- die Mordell-Vermutung, von \Person{Faltings} 1983 mit Methoden bewiesen, die Mordell nicht hätte lesen können (Arakelov-Theorie, Tate-Vermutung, Schafarewitsch-Vermutung), und von \Person{Vojta} 1991 nochmals, über diophantische Approximation.}
\end{minipage}
\end{center}

\schritt{Harnack}
\begin{center}
\begin{minipage}{0.9\textwidth}
\kursiv{\Person{Axel Harnack} bewies 1876 in Leipzig, dass eine reelle ebene Kurve vom Grad $d$ höchstens $\binom{d-1}{2} + 1$ Zusammenhangskomponenten hat, und konstruierte Kurven, die die Schranke erreichen. \textbf{Welche Anordnungen der Komponenten vorkommen, ist Hilberts 16.\ Problem} (1900) --- für Grad $6$ gelöst (Gudkov 1969), für Grad $8$ bis auf wenige Fälle (Viro, Orevkov), ab Grad $9$ offen. Das Geschlecht $g$ beschränkt die reelle Gestalt, bestimmt sie aber nicht: Der reelle Bruchfall dieses Kapitels ist ein Forschungsgebiet.}
\end{minipage}
\end{center}
\end{geschichte}

\begin{forschung}
\boxtitelp{Kurven heute}
\begin{center}
{\small
\renewcommand{\arraystretch}{1.45}
\begin{tabular}{@{}P{3.0cm}P{5.2cm}@{}}
\toprule
\textbf{Rationale Punkte finden} & Chabauty--Coleman und seine Verallgemeinerung (Kim 2005, Balakrishnan--Dogra 2018): für $g\ge2$ nicht nur Endlichkeit, sondern die Punkte selbst --- 2019 die \enquote{verfluchte Kurve} $X_s(13)$ vollständig gelöst.\\
\textbf{Gleichmäßigkeit} & Dimitrov--Gao--Habegger 2021: Die Zahl der rationalen Punkte einer Kurve vom Geschlecht $g$ ist beschränkt durch eine Konstante, die nur von $g$ und dem Rang abhängt --- der erste Schritt zur gleichmäßigen Mordell-Vermutung.\\
\textbf{Modulkurven} & Die Kurven, deren Punkte elliptische Kurven mit Zusatzstruktur klassifizieren; ihr Geschlecht entscheidet, welche Torsion über $\mathbb Q$ vorkommt (Mazur 1977) --- und ihre rationalen Punkte sind das Programm hinter Wiles.\\
\bottomrule
\end{tabular}}
\end{center}

\schritt{Zur zweiten Zeile}
\begin{center}
\begin{minipage}{0.88\textwidth}
\kursiv{\textbf{Faltings sagt \enquote{endlich}, aber nicht wie viele}; die Frage, ob es eine Schranke gibt, die nur vom Geschlecht abhängt --- gleichmäßig über alle Kurven --- war seit 1983 offen und galt als unerreichbar. \Person{Dimitrov}, \Person{Gao} und \Person{Habegger} bewiesen 2021 eine Schranke der Form $c(g)^{1 + r}$ mit dem Rang $r$ der Jacobischen; \Person{Kühne} entfernte 2021 die Abhängigkeit vom Rang für den Fall der Torsionspunkte. Die Methoden --- Höhen, o-minimale Geometrie, Pila--Zannier --- verbinden die Zahlentheorie mit der Modelltheorie, Band 10. Für den Leitfall ist das fern; aber die Frage \enquote{wie viele Modelle passen zu den Daten}, die der Band mit Grad und Bézout beantwortet hat, ist dieselbe Frage auf der anderen Seite des Faches.}
\end{minipage}
\end{center}
\end{forschung}

\begin{merksatz}
\boxtitel{Was Sie jetzt können}
\begin{center}
\renewcommand{\arraystretch}{1.5}
\begin{tabular}{@{}P{3.0cm}P{5.6cm}@{}}
\toprule
\textbf{Definition 12.1} & $HP_C = ed + 1 - p_a$\\
\textbf{Satz 12.2} & $p_a = \binom{d-1}{2}$ eben; $g = p_a - \delta$\\
\textbf{Satz 12.3} & $g = 0$ rational, $g = 1$ elliptisch, $g\ge2$ endlich viele Punkte\\
\textbf{Fermat} & Geschlecht $\binom{n-1}{2}$; Faltings gibt endlich, Wiles gibt null\\
\textbf{Der Bruchfall} & arithmetisches $\ne$ geometrisches Geschlecht; reelle Gestalt\\
\bottomrule
\end{tabular}
\end{center}

\begin{center}
\begin{minipage}{0.88\textwidth}
\kursiv{Teil IV ist zu Ende: Gruppe, Parametrisierung, Geschlecht. Teil V wendet den Band an --- Kapitel 13 auf die Statistik und den Leitfall in voller Allgemeinheit, Kapitel 14 auf Roboter und Verschlüsselung.}
\end{minipage}
\end{center}
\end{merksatz}
